\section{Reproduction and repository integration}
\label{sec:reproduction}

\subsection{Archive organization}

The archive contains this PDF, the complete \LaTeX{} source and generated
figures, a portable Python source snapshot, the additive integration patch,
raw audit and benchmark JSON, the frozen external corpus, two explicit
counterexamples, and a manifest of source identities and file hashes.
The top-level \code{README.md} gives the short entry points;
\code{INTEGRATION.md} describes the API contracts and the intended review.
The baseline commit is recorded in full so the work can be compared with
later changes to the repository.

The runnable snapshot is under \code{code/fast}. It preserves the entire
\code{fastunknot} package because the maintained package initializer imports
the existing recognizer infrastructure. Copying only a few local modules
would not preserve ordinary package imports. The snapshot includes only the
selected regression files and the fixture dependencies they require; it does
not carry unrelated historical benchmark directories. The original MIT-0
license and the internal attribution files are retained.

The new functionality consists of six modules: the weighted producer and
independent verifier, the geometric weight construction, the component
census and verifier, and the canonical disc-count wrapper. Two new test files
and four audit or timing drivers accompany them. These additions do not
modify the maintained recognizer's defaults. The integration patch is
intended to be applied at the repository root after checking the pinned
dependency identities; the full snapshot is for reproduction, not an
instruction to overwrite a newer checkout.

\subsection{Commands and optional dependencies}

Python 3.10 or newer is required by the maintained package. The weighted and
native finite-geometry code uses the standard library. Regina is needed for
the external normal-surface audit and the Regina-specific test methods;
those unit methods are reported as skipped if it is absent. Matplotlib is
needed only to regenerate the figures. A \LaTeX{} installation with the
packages named in the main source is needed only to rebuild the PDF.

From the unpacked archive, use:
\begin{quote}\small
\begin{minipage}{\linewidth}
\begin{verbatim}
python3 reproduce.py tests
python3 reproduce.py audit
python3 reproduce.py regina
python3 reproduce.py benchmark
python3 reproduce.py figures
python3 reproduce.py paper
python3 reproduce.py verify-manifest
\end{verbatim}
\end{minipage}
\end{quote}

The benchmark command runs the recorded full experiment sizes, so it takes
longer than the finite audits. The launcher uses the current Python
interpreter and supplies the import paths explicitly. Recomputed results go
to \code{reproduced/}, preserving the archived measurements in
\code{results/}. The Regina command reads the frozen corpus instead of
reenumerating vertex solutions, reducing sensitivity to enumeration order
in another Regina release. The generator can also construct a new corpus
when that is the intended experiment.

Tables and plots are generated from raw results by
\code{article/make_figures.py}; the corresponding CSV is included. A paper
rebuild uses the archived figures and tables unless they are explicitly
regenerated. The main \LaTeX{} file uses a self-contained bibliography, so
BibTeX is unnecessary. The manifest check compares bytes; it does not claim
that independently rerun timing JSON or PDF timestamps must match byte for
byte.

\subsection{Minimal native use}

The following example runs entirely on the included finite fixture. Its
expected count is two even though the total boundary class is zero.
\begin{quote}\small
\begin{minipage}{\linewidth}
\begin{verbatim}
from normal_orbit_research.fixtures import layered_torus
from fastunknot.normal_surface_components import normal_component_census
from fastunknot.normal_component_verify import (
    verify_normal_component_certificate,
)

tri, meridian = layered_torus(4)
vector = [[2 * value for value in row] for row in meridian]
answer = normal_component_census(tri, vector,
                                 record_certificate=True)
assert answer['compressing_disk_components'] == 2
assert verify_normal_component_certificate(
    tri, vector, answer['certificate'])
\end{verbatim}
\end{minipage}
\end{quote}

The separate executable example also creates a large input with coordinate
gcd one, verifies its canonical reduction, and saves both proof objects.
This gives an integration reviewer an immediately inspectable positive case,
an independent replay path, and explicit evidence of the source contract.

\subsection{Acceptance criteria for the next integration}

A reasonable first merge criterion is preservation of the old normal and
interval APIs, successful replay of the saved corpus, rejection of the
recorded source and histogram forgeries, and a caller-facing distinction
between a complete query and an inconclusive work limit. A later recognizer
integration should additionally specify the provenance witness and the
candidate-generation contract before using a positive disc count as a knot
decision. The benchmark scripts should remain usable as those later stages
are introduced, with separate measurements for discovery, transport,
reconstruction and verification.

The concrete result available now is a compressed, checked operation that
was missing from the native backend, together with a smaller exact input
for its essential-disc count. The next mathematical milestone is a bound on
the complete search or hierarchy that calls such operations, with every
representation change and bit cost charged.
